\section{What was tried to obtain a smaller exponent order}
\label{sec:obstructions}

The objective of the investigation was stronger than another small
coefficient improvement. Several routes were examined specifically because
they might remove an entire logarithmic factor. Their status is recorded
here to distinguish a missing proof from a completed improvement.

\subsection{Lower initial precision and conserved incidence degree}
The compressed dense reduction has wire budget $M=O(n\log\log n)$ but
initial counting precision only $n+1$. This suggests exploiting
$P_0\ll M$. However the generic deformation has degree $M(2L-1)$;
Hermite recovery introduces precision of leading size $6MD/u$ regardless
of that smaller initial precision. Merely replacing $P_0$ by $n+1$
does not change the dominant term.

There is a genuine algebraic alternative: if a degree bound $K=O(n)$
were retained through the compressed deformation and every recursive
entry interface, the local mode entropy could scale as
$\sqrt{MK/u}$ instead of $M/\sqrt u$. With $M=O(n\log\log n)$ this
would target $n\sqrt{\log\log n/\log n}$. The reviewed corpus already
states this conditional route. Its recent degree-flow lemma shows that
incoming mode extraction removes exactly the degree allowance transferred
to outgoing linear forms. That conservation identity is useful but does
not close the algorithmic representation.

Binary ranks of low-weight subsets do not automatically give low Boolean
polynomial degree. Carrying the original incidence polynomial preserves
degree but loses the compressed wire count; carrying a whole grouped
coefficient as one implicit tensor can make entry extraction expensive.
For example, a coefficient of a product of linear forms can be a permanent.
That example is an interface warning, not a hardness proof for the actual
Cauchy tensors. We did not construct the missing uniform compiler.

\subsection{Precision peeling and ordinary tree improvements}
An exact mode contribution with total valuation $V$ needs only residual
precision $p_t-V$. This is already an effective kernel technique in the
corpus. It does not alone justify restarting every interpolation scale.
Budget layers with $p_t-V=p_t/(\log s)^B$, for fixed $B$, can retain
the leading weighted-partition entropy while losing only
$O(\log\log s)$ in the binary logarithm of the precision. Resetting the
cumulative packet exponent by $\Theta(\log s)$ would require a much
larger reduction. This is a limitation of this budget-envelope argument,
not a universal lower bound on precision-aware algorithms.

The new affine-plane construction also lives inside a model with a known
order obstruction. Full boundary tables on a complete mode graph must
include a substantial cut in the Euclidean entropy relaxation. Root packing
simultaneously limits $u_0$ to $O(\log s/\log r)$. These statements,
already proved in the incoming packets, prevent tree selection and a faster
matrix exponent alone from proving $o(s/\sqrt{\log s})$ in that
certificate model \cite{rectangular,portfolio}. Our improvement is
consistent with that obstruction: it decreases the coefficient while
retaining its order.

\subsection{A new scoped obstruction for direct odd-prime replacement}
Changing from two-adic to $p$-adic arithmetic initially looks attractive:
an exact count needs fewer base-$p$ digits, and a growing characteristic
might appear to save a factor of $\sqrt{\log p}$. The binary locality
mechanism explains why that inference is not valid for a direct replacement.

\begin{lemma}[Leading-digit locality]\label{lem:digit}
Let $p$ be prime and let $A,B:\F_p\to\mathcal C$ be row and column
labels. If
\[
 A(x)=B(x+c)\quad\text{for all }x\in\F_p, c\in C_0,
\]
and $C_0$ contains two distinct elements, then both labels are constant.
\end{lemma}
\begin{proof}
For $c\ne d$ in $C_0$ one has $B(x+c)=B(x+d)$ for every $x$.
Thus $c-d$ is a nonzero additive period of $B$. A nonzero element
generates the additive group of the prime field, so $B$ is constant;
the displayed identity makes $A$ constant as well.
\end{proof}

In the binary construction every integer of valuation $i$ flips bit $i$,
and there is just one nonzero leading residue. For odd $p$, asking one
fixed nontrivial digit grouping to localize every nonzero leading residue
falls under Lemma~\ref{lem:digit}. Node-dependent groupings and other
deformations are outside this statement.

One can try to retain only leading digit $1$. The following exact
calculation shows the resulting precision penalty for the ordinary linear
taper. Work on the complete grid $\{0,\ldots,p^{a+1}-1\}$, excluding
zero from the interpolation nodes. Retain a node $q$ of valuation $j<u$
only when its first nonzero base-$p$ digit is $1$, and give it multiplicity
$m_j=h(u-j)$; all other multiplicities are zero. For an active node $t$
of valuation $i$ set $D=p^{a-i}$ and $H=u-i$.

\begin{proposition}[Loss of tapered cancellation]\label{prop:oddprime}
For the Hermite prefactor
\[
 E_t(0)=\prod_{q\ne0,t}\left(\frac{-q}{t-q}\right)^{m(q)},
\]
its exact valuation is
\begin{align}
 v_p(E_t(0))
 &=hD\sum_{j=1}^{a-i}\frac{j}{p^j}
       \bigl((H-j)_+-(p-1)H\bigr)\notag\\
 &=-hD\left((p-2)H\sum_{j=1}^{a-i}\frac{j}{p^j}
       +\sum_{j=1}^{a-i}\frac{j\min(j,H)}{p^j}\right).
 \label{eq:odd-loss}
\end{align}
For fixed odd $p$, at $i=0$ this has a term of order $hDu$;
the binary case cancels that term exactly.
\end{proposition}
\begin{proof}
Nodes of valuation less than $i$ have equal valuations in numerator and
denominator and contribute zero. For $j\ge1$, the grid contains
$D/p^j$ nodes of valuation $i+j$ with selected leading digit; their
multiplicity is $h(H-j)_+$. It contains $(p-1)D/p^j$ nodes with
$v_p(t-q)=i+j$; they all have valuation $i$ and the same leading digit
as $t$, hence multiplicity $hH$. Subtract these contributions and sum
the valuation level $j$. The second equality uses
$(H-j)_+=H-\min(j,H)$.
\end{proof}

Covering a deformation of degree $\Theta(Mp^{a+1})$ requires
$h=\Theta_p(M/u)$, since
$\sum_{j<u}p^{a-j}h(u-j)=\Theta_p(hp^au)$.
Thus for fixed odd $p$ the new term in \eqref{eq:odd-loss} has size
$\Theta_p(MD)$, losing the helpful $1/u$ factor. This calculation
does not rule out all nonbinary approaches. It rules out the tempting
argument that changing the radix, while keeping this locality and
one-digit taper, automatically improves the asymptotic order.

\subsection{Unbalanced binary trees do not evade the flag witness at the used arities}
The incoming flag construction proves optimality among balanced binary
separator trees in its Euclidean relaxation. We checked the more general
tree shapes at the arities $4,8,16$ used by its best polynomial-space
schedule.
Give directed pairs incident with a distinguished vertex weight $1$ and
all other directed pairs weight $2/3$. Let $P(n)$ be the smallest possible
maximum root-to-leaf cut cost on $n$ ordinary vertices, and $S(n)$ the
corresponding value with one distinguished vertex. Exact recurrences are
\begin{align*}
 P(1)&=S(1)=0,\\
 P(n)&=\min_{1\le a<n}\left\{\frac43a(n-a)
                      +\max(P(a),P(n-a))\right\},\\
 S(n)&=\min_{1\le a<n}\left\{\frac23(n-a)(2a+1)
                      +\max(S(a),P(n-a))\right\}.
\end{align*}
In the last recurrence the distinguished child has $a$ vertices.
All labels of the same type are interchangeable, so every binary tree
shape is covered; the recurrences are proved by its first split.
Exact rational evaluation gives
\[
 S(4)=26/3,\qquad S(8)=98/3,\qquad S(16)=370/3.
\]
These equal the squared witness norms in the existing flag certificate.
After normalization, the optimal local coefficients at these arities
remain $r^{-1}\sqrt{(r-1)(4r+10)/9}$ even when unbalanced trees are
allowed. This is a finite model-specific strengthening. It neither proves
the same formula for every $r$ nor lower-bounds arbitrary contraction
algorithms; indeed a naive extension of the formula to every integer
arity is false.
